\clearpage\begingroup
\SourceLeaf{143b}\HistoricalHead{138}{L}{PHYSIK.\quad NACHLASS.}
\makeatletter\tagsleft@true\let\veqno\@@leqno\makeatother
\begin{center}\textbf{ERLÄUTERUNGEN.}\end{center}
\noindent\hspace*{\parindent}Die \textsc{Gauss}sche Handschrift enthält einen Teil der Rechnungen für das von ihm angegebene Doppelobjektiv (vergl. Werke V, S. 504 ff.). \textsc{Gauss} stützt sich hierbei ganz auf \textsc{Euler}s \textit{Dioptricae pars prima continens librum primum, de explicatione principiorum ex quibus constructio tam telescopiorum quam microscopiorum est petenda}, Petropoli 1769; wir zitieren dieses Werk im folgenden nach der neuen Ausgabe: \textsc{Leonhardi Euleri} Opera omnia, series III, vol. 3, 1911.

\textsc{Gauss} knüpft an die drei letzten Formeln auf S. 34 der Dioptrik \textsc{Euler}s an, die folgendes Problem lösen: Bei gegebenem Objektabstand $a$ und gegebenem Bildabstand $\alpha$ gibt es unendlich viele, als unendlich dünn betrachtete Einzellinsen, die dann das Objekt in dem Bildpunkt abbilden; allerdings ist die sphärische Aberration (\textsc{Euler}s »spatium diffusionis«) je nach Wahl der Linse verschieden groß, und einer bestimmten Linse entspricht das Minimum derselben. In \textsc{Gauss}scher Schreibweise lauten die genannten \textsc{Euler}schen Formeln für die Radien dieser Linsen folgendermaßen:
\begin{equation*}\tag*{($\alpha$)}
\left\{\begin{aligned}
\frac1\rho&=\frac{a\alpha}{r\alpha+sa\pm\tau(a+\alpha)\sqrt{l-1}},\\[5pt]
\frac1\sigma&=\frac{a\alpha}{ra+s\alpha\mp\tau(a+\alpha)\sqrt{l-1}},
\end{aligned}\right.
\end{equation*}
\noindent und die sphärische Aberration wird gleich
\begin{equation*}\tag*{($\beta$)}
P\alpha^2x^2=\mu\alpha^2x^2\left(\frac1a+\frac1\alpha\right)
\left[l\left(\frac1a+\frac1\alpha\right)^2+\frac\nu{a\alpha}\right];
\end{equation*}
\noindent $x$ ist dabei die Apertur der (ersten) Linse und $r,s,\tau,\mu,\nu$ haben die folgende Bedeutung, wenn $(m+1)$ der Brechungsexponent der (ersten) Linse für rotes Licht ist:
\begin{equation*}\tag*{($\gamma$)}
\left\{\begin{aligned}
\mu&=\frac{(m+1)(4m-3)}{8m^2(m+3)},\\[4pt]
\nu&=\frac{4m^2}{4m+3},\\[4pt]
r&=\frac1{2m}+\frac1{m+3}-1,\\[4pt]
s&=1+\frac1{2m}-\frac1{m+3},\\[4pt]
\tau&=\frac{(m+1)\sqrt{4m+3}}{2m(m+3)}.
\end{aligned}\right.
\end{equation*}
\noindent Ferner ist $l$ ein Parameter $\geqq1$; setzt man ihn gleich $1$, so erhält man diejenige Linse, für die die sphärische Aberration ein Minimum wird. Denkt man sich eine zweite Linse, die das Bild der ersten zum Objekte hat (Objektabstand $b$) und dieses im Bildabstande $\beta$ abbildet, so lauten dafür die entsprechenden
\clearpage\SourceLeaf{144}\HistoricalHead{139}{R}{ACHROMATISCHE DOPPELOBJECTIVE, OHNE RÜCKSICHT AUF DICKE UND ABSTAND.}
\noindent Formeln:
\begin{equation*}\tag*{($\delta$)}
\left\{\begin{aligned}
\frac1{\rho'}&=\frac{b\beta}{r'\beta+s'b\pm\tau'(b+\beta)\sqrt{l'-1}},\\[5pt]
\frac1{\sigma'}&=\frac{b\beta}{r'b+s'\beta\mp\tau'(b+\beta)\sqrt{l'-1}},\\[5pt]
\text{Aberration}&=Q\beta^2x'^2=\mu'\beta^2x'^2\left(\frac1b+\frac1\beta\right)\\
&\qquad\left[l'\left(\frac1b+\frac1\beta\right)^2+\frac{\nu'}{b\beta}\right].
\end{aligned}\right.
\end{equation*}
\noindent\hspace*{\parindent}Die gestrichelten Größen $\mu',\nu',r',s',\tau'$ sind nach ($\gamma$) mit $(m'+1)$, dem Brechungsexponenten der zweiten Linse für rotes Licht zu bilden; $x'$ ist die Apertur der zweiten Linse.

Steht die zweite Linse unmittelbar hinter der ersten, so ist offenbar:
\begin{equation*}\tag*{($\varepsilon$)}b=-\alpha,\end{equation*}
\noindent und die Gesamtaberration beider Linsen ist nach \textsc{Euler} (\textit{Dioptrica}, S. 65 und 75):
\[\text{Gesamtaberration}=P+Q,\]
\noindent wo $P$ und $Q$ die aus ($\beta$) und der letzten Gleichung ($\delta$) ersichtlichen Abkürzungen sind.

Wir führen nun die speziellen Bedingungen und Bezeichnungen der \textsc{Gauss}schen Handschrift ein. Nach \textsc{Gauss} ist:
\[a=\infty,\quad\beta=\omega,\quad b=-\alpha=p;\]
\noindent setzt man diese Werte in die beiden Gleichungen ($\alpha$) ein, so folgt für $\rho$ und $\sigma$:
\begin{equation*}\tag*{($\zeta$)}
\left\{\begin{aligned}
\rho&=\frac sp\mp\frac{\tau\sqrt{l-1}}p,\\[4pt]
\sigma&=\frac rp\mp\frac{\tau\sqrt{l-1}}p,
\end{aligned}\right.
\end{equation*}
\noindent also:
\[\rho+\sigma=\frac{s+r}p,\]
\noindent oder nach Benutzung von ($\gamma$):
\begin{equation*}\tag*{($\eta$)}\rho+\sigma=\frac1{mp},\end{equation*}
\noindent was mit der \textsc{Gauss}schen Gleichung [1] übereinstimmt.

Ebenso folgt aus ($\zeta$) und ($\gamma$) die \textsc{Gauss}sche Gleichung [2]:
\begin{equation*}\tag*{($\vartheta$)}
\begin{aligned}
\rho-\sigma&=\frac{s-r}p\pm\frac{2\tau\sqrt{l-1}}p\\
&=\frac{2m+4}{(m+3)p}\pm\frac{(m+1)\sqrt{4m+3}\sqrt{l-1}}{m(m+3)p};
\end{aligned}
\end{equation*}
\noindent \textsc{Gauss} benutzt im Folgenden nur das obere Zeichen.

In derselben Weise folgen aus ($\delta$) und den in den gestrichelten Größen geschriebenen Gleichungen ($\gamma$) die Formeln:
\[\begin{aligned}
\rho'+\sigma'&=\frac1{m'}\left(\frac1\omega-\frac1p\right),\\[4pt]
\rho'-\sigma'&=\frac{2m'+4}{m'+3}\left(\frac1\omega+\frac1p\right)\\
&\quad+\frac{(m'+1)\sqrt{4m'+3}}{m'(m'+3)}\left(\frac1\omega-\frac1p\right)\sqrt{l'-1},
\end{aligned}\]
\noindent d.\,h. die \textsc{Gauss}schen Gleichungen [3] und [4], S. 135.

Die Forderung, daß die gesamte sphärische Aberration, die durch die Linsenkombination hervorgebracht ist, verschwinde, ist nach dem Vorhergehenden offenbar $P+Q=0$, d.\,h., da für unsere
\par\HistoricalFoot{}{18*}
\clearpage\SourceLeaf{145}\HistoricalHead{140}{L}{PHYSIK.\quad NACHLASS.}
\noindent Bedingungen
\[\begin{aligned}
P&=\frac\mu{p^3}l,\\[4pt]
Q&=\mu'\left\{\frac1\omega-\frac1p\right\}\left\{l'\left(\frac1\omega-\frac1p\right)^2-\frac{\nu'}{p\omega}\right\}
\end{aligned}\]
\noindent ist, so folgt die Gleichung:
\begin{equation*}\tag*{($\iota$)}
\frac{\mu l}{p^3}+\mu'\left\{\frac1\omega-\frac1p\right\}\left\{l'\left(\frac1\omega-\frac1p\right)^2-\frac{\nu'}{p\omega}\right\}=0,
\end{equation*}
\noindent die durch Einsetzen der Werte $\mu,\mu',\nu,\nu'$ u.s.w. nach Gleichung ($\gamma$) übergeht in die \textsc{Gauss}sche Gleichung [5]:
\begin{equation*}\tag*{($\kappa$)}
\begin{aligned}
&\frac{(m+1)(4m+3)}{m^2(m+3)}\frac l{p^3}
+\frac{(m'+1)(4m'+3)l'}{m'^2(m'+3)}\left(\frac1\omega-\frac1p\right)^3\\
&\qquad-\frac{4(m'+1)}{m'+3}\frac1{\omega p}\left(\frac1\omega-\frac1p\right)=0.
\end{aligned}
\end{equation*}
\noindent Nun ist nach der gewöhnlichen Linsenformel die reziproke Brennweite einer Linse gleich $\dfrac1a+\dfrac1\alpha$, und man hat weiter
\[\frac1a+\frac1\alpha=(n-1)(\rho+\sigma),\]
\noindent wenn man für einen Augenblick den Brechungsexponenten für eine beliebige Farbe mit $n$ bezeichnet. Da nach \textsc{Gauss} die Brennweite (\textsc{Gauss} schreibt »erste Vereinigungsweite«) der ersten Linse für rotes Licht $p$, für violettes $P$ genannt wird, und in seiner Bezeichnungsweise die entsprechenden Brechungsexponenten $m+1$ bezw.\ $(M+1)$ sind, so folgt durch zweimalige Anwendung der letzten Formel:
\[\frac1p=m(\rho+\sigma),\qquad\frac1P=M(\rho+\sigma),\]
\noindent also:
\begin{equation*}\tag*{($\lambda$)}\frac1{mp}=\frac1{MP}.\end{equation*}
\noindent Für die zweite Linse gilt ebenso, wenn wir zunächst die alten Bezeichnungen $b$ und $\beta$ wiederaufnehmen.
\[\frac1b+\frac1\beta=(n'-1)(\rho'+\sigma'),\]
\noindent wobei $n'$ der Brechungsexponent für eine beliebige Farbe sei. Für rotes Licht ist nach dem Früheren $b=-p$, für violettes dagegen $b=-P$ zu nehmen. Da ferner die \textsc{Gauss}schen Brechungsexponenten bezw.\ $(m'+1)$ und $(M'+1)$ sind, so ergibt die zweimalige Anwendung dieser Gleichung:
\[\begin{aligned}
-\frac1p+\frac1{\beta_{\mathrm{rot}}}&=m'(\rho'+\sigma'),\\[4pt]
-\frac1P+\frac1{\beta_{\mathrm{viol.}}}&=M'(\rho'+\sigma').
\end{aligned}\]
\noindent Jetzt verlangt die Forderung, daß die Linsenkombination \textls[160]{achromatisch} sein soll (wenigstens für die sog.\ Nullstrahlen), daß $\beta_{\mathrm{rot}}=\beta_{\mathrm{viol.}}$, welchen gemeinsamen Wert wir nach \textsc{Gauss} jetzt mit $\omega$ bezeichnen. Dann gehen die beiden letzten Gleichungen über in die »Bedingung der Achromasie«:
\[\left(\frac1\omega-\frac1P\right)=\frac{M'}{m'}\left(\frac1\omega-\frac1p\right)\]
\clearpage\SourceLeaf{146}\HistoricalHead{141}{R}{ACHROMATISCHE DOPPELOBJECTIVE, OHNE RÜCKSICHT AUF DICKE UND ABSTAND.}
\noindent die vermöge ($\lambda$) in einer der folgenden 4 Formen geschrieben werden kann:
\begin{equation*}\tag*{($\mu$)}
\left\{\begin{aligned}
\frac1\omega&=\frac{M'}{m'}\left(\frac1\omega-\frac1p\right)+\frac M{mp},\\[4pt]
\left(\frac Mm-\frac{M'}{m'}\right)\frac1p&=\left(1-\frac{M'}{m'}\right)\frac1\omega,\\[4pt]
\frac p\omega&=\frac{M'm-Mm'}{m(M'-m')},\\[4pt]
1-\frac p\omega&=\frac{m'(M-m)}{m(M'-m')}.
\end{aligned}\right.
\end{equation*}
\noindent Bei vorgeschriebener Brennweite des Doppelobjektivs bestimmen die letzten Formeln $p$, und dieses nach dem Vorhergehenden $(\rho+\sigma)$ und $(\rho'+\sigma')$, so daß man sich die Linsen durch $(\rho+\sigma)$ und $(\rho'+\sigma')$ gegeben denken kann. Nach Division von [1] in [2], sowie von [3] in [4] erhält man dann unter Benutzung der \textsc{Gauss}schen Abkürzungen $e$ und $e'$ nach den Gleichungen [9] und [11], S. 136:
\begin{equation*}\tag*{($\nu$)}
\left\{\begin{aligned}
\frac{\rho-\sigma}{\rho+\sigma}&=\frac{2m(m+2)}{m+3}+\frac{(m+1)\sqrt{4m+3}}{m+3}\sqrt{l-1}\\
&=e+\frac{(m+1)\sqrt{4m+3}}{m+3}\sqrt{l-1},\\[4pt]
\frac{\rho'-\sigma'}{\rho'+\sigma'}&=e'+\frac{(m'+1)\sqrt{4m'+3}}{m'+3}\sqrt{l'-1}.
\end{aligned}\right.
\end{equation*}
\noindent Führt man ferner $\dfrac p\omega$ aus Gleichung ($\mu$) in die Bedingung der sphärischen Korrektion, d.\,h. in die \textsc{Gauss}sche Gleichung [5] ein, so nimmt diese die Gestalt an:
\[\begin{aligned}
&\frac{(m+1)(4m+3)}{m^2(m+3)}l
-\frac{(m'+1)(4m'+3)}{m'^2(m'+3)}\left(\frac{m'}m\right)^3\left(\frac{M-m}{M'-m'}\right)^3l'\\
&\qquad+\frac{4(m'+1)}{m'+3}\mathbin{.}\frac{M'm-Mm'}{m(M'-m')}\mathbin{.}\frac{m'(M-m)}{m(M'-m')}=0,
\end{aligned}\]
\noindent oder:
\begin{equation*}\tag*{($\xi$)}
\begin{aligned}
&\frac{m(m+1)(4m+3)}{(m+3)(M-m)^3}l
-\frac{m'(m'+1)(4m'+3)}{(m'+3)(M'-m')^3}l'\\
&\qquad+\frac{4mm'(M'm-Mm')}{(M'-m')^2(M-m)^2}=0,
\end{aligned}
\end{equation*}
\noindent oder endlich, unter Verwendung der \textsc{Gauss}schen Abkürzungen in den Gleichungen [6], [7] und [8] eine der drei folgenden Formen
\begin{equation*}\tag*{($\pi$)}
\left\{\begin{aligned}
al-bl'+c&=0,\\[3pt]
a(l-1)-b(l'-1)+(a-b+c)&=0,\\[3pt]
\frac a{a-b+c}(l-1)-\frac b{a-b+c}(l'-1)+1&=0.
\end{aligned}\right.
\end{equation*}
\noindent Man zeigt nun leicht, daß, wenn das erste Glas Kron, das zweite Flint ist, $a,b,c$ größer als Null sind, sowie daß $a>b$. Da $l>1$ ist, so ist also der Ausdruck $\dfrac a{a-b+c}(l-1)$ größer als Null, und man kann setzen:
\[\frac a{a-b+c}(l-1)=\operatorname{tang}^2\psi.\]
\noindent Damit schreibt sich die letzte Gleichung ($\pi$), d.\,h. die Bedingung für sphärische Korrektur:
\[\frac b{a-b+c}(l'-1)=\sec^2\psi.\]
\noindent Einsetzen in die Gleichungen ($\nu$) unter Benutzung der \textsc{Gauss}schen Abkürzungen $f$ und $f'$ liefert die
\clearpage\SourceLeaf{147}\HistoricalHead{142}{L}{PHYSIK.\quad NACHLASS.}
\noindent Gleichungen:
\begin{equation*}\tag*{($\rho$)}
\left\{\begin{aligned}
\frac{\rho-\sigma}{\rho+\sigma}&=e+f\operatorname{tang}\psi,\\[4pt]
\frac{\rho'-\sigma'}{\rho'+\sigma'}&=e'+f'\sec\psi,
\end{aligned}\right.
\end{equation*}
\noindent d.\,h. die \textsc{Gauss}schen Gleichungen [Ia] und [Ib], S. 136. Beliebige Wahl des Parameters $\psi$ liefert also alle Doppelobjektive, die in der Achse chromatisch und für rotes Licht sphärisch korrigiert sind.

Vertauscht man $m$ mit $M$, $m'$ mit $M'$, so erhält man die zweiteiligen Linsen, die in der Achse chromatisch und für violettes Licht sphärisch korrigiert sind. Das liefert die \textsc{Gauss}schen Gleichungen [IIa] und [IIb], S. 136.

Die Behandlung der zweiten Lösungsform von \textsc{Gauss} ist ganz analog. — Es ist zu bemerken, daß man, wenn nur die Brennweite vorgeschrieben wird, nicht zu dem Doppelobjektiv gelangt, das \textsc{Gauss} a.a.O. angegeben hat. Denn hier sind Dicke der Linsen und ihr Abstand vernachlässigt. Indessen dient diese Rechnung, wie \textsc{Gauss} selbst hervorhebt, als Ausgangspunkt, und die hier erhaltenen Näherungswerte werden für die genauere Rechnung benutzt. Diese findet sich in der Handschrift nicht.

Die erste öffentliche Erwähnung der \textsc{Gauss}schen Formeln I. und II. findet sich übrigens in einem Briefe von \textsc{Gauss} an \textsc{Schumacher} vom 25. Juni 1810, \textit{Briefwechsel zwischen Gauss und Schumacher} I, Altona 1860, S. 45 (siehe weiter unten S. 146), in dem \textsc{Gauss} diese Form der Behandlung als ihm eigentümlich bezeichnet. Mittels dieser Formeln hatte er damals ein Doppelobjektiv für \textsc{Repsold} in Hamburg berechnet, über dessen Ausfall die folgenden Briefe des Briefwechsels \textsc{Gauss}–\textsc{Schumacher} Auskunft geben.

Der Inhalt dieser Erläuterung beruht im Wesentlichen auf einer schriftlichen Mitteilung von H. \textsc{Boegehold} in Jena.
\begin{flushright}\textsc{Schaefer.}\end{flushright}
\begin{center}\rule{.25\textwidth}{.4pt}\end{center}
\clearpage\endgroup
